\section{おわりに}
本プロジェクトでは，幼児教育での活用を目的として，複数の楽器が同期して演奏し，あわせて演奏に
同期して動作する人形による視覚的演出を行う分散型の自動演奏システムを開発した．このうち報告者は，
トロンボーンおよびカスタネットの音源の設計と実装を担当した．

トロンボーンの音源では，オシレータによる波形にローパスフィルタ，ディレイ，および音量エンベロープを
適用することで，管楽器に近い音色を合成した．カスタネットの音源では，オシレータではなくノイズを
音源とし，これを帯域通過フィルタと音量エンベロープに通すことで，音程を持たない打撃音を再現した．

評価の結果，客観評価においては，担当したいずれの楽器も理論周波数との誤差が$\pm1.0$\,\%以内という
基準を満たし，実曲の演奏においても発音成功率100\,\%を達成した．一方，主観評価では，作成音の平均
評価値が基準の4.0に達しなかった．これは，合成音の音量や，他の楽器との音量・音域の関係によって，
担当した楽器の音が相対的に埋もれたことが要因と考えられる．また，トロンボーンの音色設計においては，
生成AIを用いた倍音の抽出に誤りがあったものの，フィルタおよびエンベロープによる処理によって，
管楽器に近い音色を得ることができた．

今後の課題として，音源の音量や音量バランスの調整，および実際の楽器の倍音構造を正確に抽出して
波形に反映することが挙げられる．これらの改善により，客観評価に加えて主観評価においても基準を
満たす音源の実現が期待される．本システムは，音と視覚的演出を組み合わせることで，幼児のリズム
教育に活用できる自動演奏システムとしての可能性を示すものである．